\scanpage{22}
On définit aussi
\[
T^n(E)=\operatorname{cl}_n(\operatorname{Hom}(I_n,E))\in k_n(G)
\]
\margin{$(\unsure{E\in k(G)})$}
\[
T^*(E)=\sum T^n(E)\in1+\widehat{k_*(G)}^+
\]
on donc encore
\[
T^*(E\amalg F)=T^*(E)T^*(F)
\]
d'où
\[
T^*:k_0(G)\mathbin{\text{---}}1+\widehat{k_*(G)}^+
\]
défini de façon unique par
\[
T^*(\operatorname{cl}_0(E))=T^*(E)
\]
On veut donner une relation entre $T^*$ et $\tau^*$,
permettant de calculer $T^*$ en termes de $\tau^*$.
\[
\begin{aligned}
\operatorname{Hom}(I_n,E)&\simeq\coprod_{\substack{Q\text{ quotient}\\\text{de }I_n}}\operatorname{Mon}(Q,E)\\
&\simeq\coprod_{\substack{\deleted{\unsure{1\leq i\leq n}}\\\alpha=(\alpha_1,\ldots,\alpha_n)\in\mathbf N^n\\\sum i\alpha_i=n}}\underbrace{\coprod_{\substack{Q\text{ quotient}\\\text{de }I_n\text{ de}\\\text{type }\alpha}}\operatorname{Mon}(Q,E)}_{\varphi_{\unsure{n,i|\alpha,i}}(\tau^i(E))}
\end{aligned}
\]
\deleted{donc
\[
T^n(E)=\sum_{1\leq i\leq n}\varphi_{n,i}(\tau^i(E))
\]
où
\[
\varphi_{n,i}:k_i(G)\mathbin{\text{---}}k_n(G)
\]
est défini par
\[
\varphi_{n,i}(\operatorname{cl}_i(E))=\operatorname{cl}_n(E\mathbin{\wedge}_{\mathfrak S_i}\mathfrak S_n)
\]
\[
\begin{aligned}
T^*(E)&=\sum_{n\geq0}T^n(E)=1+\sum_{1\leq i\leq n}\varphi_{n,i}\tau^i(E)\\
&=1+\sum_{i\geq1}\varphi_{\unsure{*i}}\tau^i(E)
\end{aligned}
\]
}
\mspage{21}
